% ==========================================================================
%  Chapter 46 — Phase 1: Constitutive State Protocol — Audit and Validation
% ==========================================================================

\chapter{Phase 1: Constitutive State Protocol}
\label{ch:phase1-constitutive-state}

This chapter documents Phase~1 of the master plan
(Chapter~\ref{ch:convergence-master-plan}).  The original scope
called for introducing a trial/committed state protocol through the
full constitutive chain.  A deep audit of the codebase revealed that
this protocol \emph{already exists} across every architectural layer,
from storage policy up to the SNES solver boundary.  This chapter
records the audit, maps the complete commit chain with its mathematical
foundations, and identifies the minor gaps that remain for future phases.

%% ------------------------------------------------------------------
\section{Phase 1 Scope Revision}
\label{sec:phase1-scope-revision}

Chapter~\ref{ch:convergence-master-plan} defined four Phase~1 deliverables:

\begin{enumerate}
  \item Define a C++20 concept \texttt{ConstitutivePolicy} requiring
        \texttt{commit}, \texttt{compute\_response}, \texttt{tangent}.
  \item Refactor \texttt{MaterialPoint} so every quadrature site stores
        committed \emph{and} trial state.
  \item Adapt elastic constitutive policies to satisfy the concept.
  \item Implement one inelastic model with \texttt{commit}/\texttt{revert}
        (Von~Mises $J_2$ plasticity).
\end{enumerate}

After a thorough audit of \texttt{src/materials/},
\texttt{src/elements/}, and \texttt{src/analysis/}, all four
deliverables were found to be \textbf{already implemented}:

\begin{center}
\begin{tabular}{clc}
\toprule
\textbf{\#} & \textbf{Deliverable} & \textbf{Status} \\
\midrule
1 & Constitutive concept hierarchy &
    \texttt{ConstitutiveRelation} +
    \texttt{ConstitutiveIntegratorConcept} \\
2 & Trial/committed state at each site &
    \texttt{TrialCommittedState<S>} +
    \texttt{ConstitutiveState<\ldots>} \\
3 & Elastic adaptation &
    \texttt{PassThroughIntegrator} +
    \texttt{ElasticRelation} \\
4 & Inelastic model with commit &
    \texttt{PlasticityRelation} +
    \texttt{EmbeddedRelationIntegrator} \\
\bottomrule
\end{tabular}
\end{center}

Consequently, Phase~1 was re-scoped to:
(a)~comprehensive audit and documentation of the existing protocol,
(b)~verification that the commit chain is complete from analysis to
storage, and
(c)~identification of gaps for future phases (notably the
absence of an explicit \texttt{revert} at the type-erased boundary).

%% ------------------------------------------------------------------
\section{Mathematical Framework}
\label{sec:phase1-mathematics}

\subsection{Trial/Committed State Protocol}

The incremental solution of a nonlinear boundary-value problem by
Newton--Raphson requires a clear distinction between \emph{converged}
(committed) and \emph{tentative} (trial) states at each constitutive
site.  Denoting by $n$ the load step and by $k$ the Newton iteration:

\begin{itemize}
  \item \textbf{Committed state} $\boldsymbol{\alpha}_n$:
        the internal variable set at the end of step~$n$, after global
        convergence has been verified.  This is the \emph{reference}
        state for step~$n+1$.
  \item \textbf{Trial state} $\boldsymbol{\alpha}_{n+1}^{(k)}$:
        the tentative internal variables computed during iteration~$k$
        of step~$n+1$.  This state must \emph{never} overwrite the
        committed state until the global solver converges.
\end{itemize}

The protocol at each constitutive site is:

\begin{enumerate}
  \item \textbf{Evaluate} (side-effect free):
        Given the trial strain $\boldsymbol{\varepsilon}_{n+1}^{(k)}$
        and committed internal variables $\boldsymbol{\alpha}_n$,
        compute the stress response and algorithmic tangent:
        \begin{align}
          \boldsymbol{\sigma}_{n+1}^{(k)}
            &= \hat{\boldsymbol{\sigma}}\!\left(
                 \boldsymbol{\varepsilon}_{n+1}^{(k)},\,
                 \boldsymbol{\alpha}_n \right), \\[3pt]
          \mathbb{C}_{n+1}^{(k)}
            &= \frac{\partial\hat{\boldsymbol{\sigma}}}
                    {\partial\boldsymbol{\varepsilon}}
               \bigg|_{
                 \boldsymbol{\varepsilon}_{n+1}^{(k)},\,
                 \boldsymbol{\alpha}_n}.
        \end{align}
        These evaluations do \emph{not} modify the internal state.

  \item \textbf{Commit} (after global convergence):
        Promote trial state to committed:
        \[
          \boldsymbol{\alpha}_{n+1} \leftarrow
            \hat{\boldsymbol{\alpha}}\!\left(
              \boldsymbol{\varepsilon}_{n+1}^{(\text{conv})},\,
              \boldsymbol{\alpha}_n \right).
        \]

  \item \textbf{Revert} (on solver divergence):
        Discard all trial computations; internal state remains at
        $\boldsymbol{\alpha}_n$.  In the current SNES-driven
        architecture, revert is implicit: if Newton does not converge,
        \texttt{commit()} is simply never called.
\end{enumerate}

\subsection{Constitutive Tangent in the Newton System}

The global Newton update at iteration $k$ solves
\begin{equation}
  \underbrace{
    \mathbf{K}_t^{(k)}
  }_{\text{tangent stiffness}}
  \,\delta\!\mathbf{u}^{(k)}
  = -\underbrace{
    \left(\mathbf{f}_{\mathrm{int}}^{(k)} - \mathbf{f}_{\mathrm{ext}}\right)
  }_{\text{residual}},
  \label{eq:newton-update}
\end{equation}
where
\[
  \left(\mathbf{K}_t\right)_{IJ}
    = \int_\Omega
        \mathbf{B}_I^T \,\mathbb{C}_{n+1}^{(k)}\, \mathbf{B}_J
      \,d\Omega
    + \underbrace{
        \int_\Omega
          \nabla N_I \cdot \boldsymbol{\sigma}_{n+1}^{(k)} \cdot \nabla N_J
        \,d\Omega
      }_{\text{geometric stiffness } K_\sigma}
\]
The geometric stiffness $K_\sigma$ is assembled only under
\texttt{TotalLagrangian} and \texttt{UpdatedLagrangian} kinematic
policies.  For \texttt{SmallStrain}, $K_\sigma = 0$.

The critical requirement is that both
$\boldsymbol{\sigma}(\boldsymbol{\varepsilon})$ and
$\mathbb{C}(\boldsymbol{\varepsilon})$ are evaluated as
\textbf{pure functions} of the trial kinematics and committed
internal state.  This guarantees that Newton iterations are
reproducible and that a diverged step leaves no corrupted state.


%% ------------------------------------------------------------------
\section{The Complete Commit Chain}
\label{sec:phase1-commit-chain}

The constitutive state protocol is implemented across six architectural
layers. This section traces the commit path from solver to storage.

\subsection{Layer 1: Storage Policy}
\label{sec:layer1-storage}

File: \texttt{src/materials/MaterialStatePolicy.hh}

The storage policy controls how state values are persisted:

\begin{description}
  \item[\texttt{ElasticState<S>}]
    Single value, no history.  \texttt{commit\_trial()} is not available
    (constrained away by SFINAE).

  \item[\texttt{TrialCommittedState<S>}]
    Stores two copies: \texttt{committed\_} and
    \texttt{trial\_}.  Provides:
    \begin{itemize}
      \item \texttt{commit\_trial()} --- copies trial to committed,
            clears the trial flag.
      \item \texttt{revert\_trial()} --- clears the trial value,
            restores the committed state.
      \item \texttt{committed\_value()} / \texttt{trial\_value()} ---
            accessors.
      \item \texttt{has\_trial\_value()} --- boolean flag.
    \end{itemize}
    This is the storage policy used by all path-dependent constitutive
    models.
\end{description}

\subsection{Layer 2: Constitutive State}
\label{sec:layer2-constitutive-state}

Files:
\texttt{src/materials/MaterialState.hh},
\texttt{src/materials/ConstitutiveState.hh}

\texttt{ConstitutiveState<StoragePolicy, StateVariableTypes...>} adds
semantic meaning above the raw storage.  It conditionally exposes
\texttt{commit\_trial()} and \texttt{revert\_trial()} only when the
underlying storage policy supports them (via C++20 \texttt{requires}
expressions).  This is the single point where the trial/committed
protocol is semantically anchored.

\subsection{Layer 3: Constitutive Integrator}
\label{sec:layer3-integrator}

File: \texttt{src/materials/ConstitutiveIntegrator.hh}

Three integrator strategies inject the algorithmic policy at the
\texttt{Material<>} boundary:

\begin{description}
  \item[\texttt{PassThroughIntegrator}]
    Elastic path.  \texttt{compute\_response()} and \texttt{tangent()}
    delegate directly to the constitutive relation.
    \texttt{commit()} is an intentional no-op.

  \item[\texttt{EmbeddedRelationIntegrator}]
    Inelastic path (compatibility bridge).  Delegates
    \texttt{compute\_response()} and \texttt{tangent()} to the
    relation (which may use an external algorithmic state).
    On \texttt{commit()}, calls:
    \begin{enumerate}
      \item \texttt{relation.update($\boldsymbol\varepsilon$)} or
            \texttt{relation.commit($\boldsymbol\alpha$,
            $\boldsymbol\varepsilon$)} --- updates internal variables.
      \item \texttt{site.constitutive\_state().commit\_trial()} ---
            promotes trial to committed.
    \end{enumerate}

  \item[\texttt{ContinuumLocalProblemIntegrator}]
    Advanced continuum path for finite-strain local problems.
    Uses a \texttt{ContinuumLocalIntegrationAlgorithm} (Newton at the
    local level).  On \texttt{commit()}, stores the converged
    algorithmic state and calls \texttt{commit\_constitutive\_state()}.
\end{description}

The key contract is captured by the concept:

\begin{verbatim}
template <typename I, typename Site>
concept ConstitutiveIntegratorConcept = requires(I& i, Site& s,
    const typename Site::KinematicT& k) {
    { i.compute_response(s, k) } -> same_as<typename Site::ConjugateT>;
    { i.tangent(s, k) }          -> same_as<typename Site::TangentT>;
    { i.commit(s, k) };
};
\end{verbatim}

\subsection{Layer 4: Type-Erased Material Handle}
\label{sec:layer4-material}

File: \texttt{src/materials/Material.hh}

\texttt{Material<Policy>} is an owning, value-semantic handle that
binds a concrete constitutive site and integrator behind a virtual
interface:

\begin{center}
\begin{tabular}{ll}
\toprule
\textbf{Method} & \textbf{Semantics} \\
\midrule
\texttt{compute\_response(k)} & Side-effect free stress evaluation \\
\texttt{tangent(k)}           & Algorithmic tangent at trial kinematics \\
\texttt{commit(k)}            & Promote trial $\to$ committed \\
\texttt{update\_state(k)}     & Cache current kinematic state \\
\texttt{clone\_owned()}       & Deep copy for heterogeneous containers \\
\bottomrule
\end{tabular}
\end{center}

The virtual dispatch is through \texttt{MaterialConcept<Policy>}, with
concrete implementation in \texttt{OwningMaterialModel<MaterialType,
IntegratorT>}.  The \texttt{commit()} virtual function delegates to
\texttt{integrator\_.commit(material\_, k)}, which triggers the full
Layer~3 commit chain.

\paragraph{Notable absence:}
The type-erased interface does \emph{not} expose \texttt{revert()}.
The \texttt{revert\_trial()} method exists at Layers~1--2 but is not
wired through the virtual table.  In the current SNES-driven
architecture this is correct, because revert is implicit (no
commit on divergence).  For future arc-length or substepping
strategies that may need explicit revert at the global level,
adding a \texttt{revert()} virtual to \texttt{MaterialConcept}
is a Phase~3 item.

\subsection{Layer 5: Finite Element}
\label{sec:layer5-element}

Files:
\texttt{src/elements/FiniteElementConcept.hh},
\texttt{src/elements/ContinuumElement.hh},
\texttt{src/elements/BeamElement.hh},
\texttt{src/elements/TimoshenkoBeamN.hh},
\texttt{src/elements/ShellElement.hh},
\texttt{src/elements/FEM\_Element.hh},
\texttt{src/elements/StructuralElement.hh}

The \texttt{FiniteElement} concept requires:

\begin{verbatim}
template <typename E>
concept FiniteElement = requires(E e, Mat K, Vec u, Vec f) {
    e.set_num_dof_in_nodes();
    e.inject_K(K);
    e.compute_internal_forces(u, f);
    e.inject_tangent_stiffness(u, K);
    e.commit_material_state(u);  // <-- the critical method
};
\end{verbatim}

Every concrete element type implements
\texttt{commit\_material\_state(Vec u\_local)}:

\begin{description}
  \item[\texttt{ContinuumElement}]
    Loops over quadrature points, computes the kinematic carrier
    (small-strain, Total Lagrangian, or Updated Lagrangian), then
    calls \texttt{material\_point.commit(kinematics)}.
    This propagates through \texttt{Material<>::commit(kin)}.

  \item[\texttt{BeamElement}, \texttt{TimoshenkoBeamN}, \texttt{ShellElement}]
    Same pattern: extract local DOFs, transform to local frame,
    compute section strain at each quadrature point, then call
    \texttt{section.commit(strain)} followed by
    \texttt{section.update\_state(strain)}.  Structural elements
    perform both commit \emph{and} state update inside
    \texttt{commit\_material\_state()}.
\end{description}

The type-erased wrappers \texttt{FEM\_Element} and
\texttt{StructuralElement} forward \texttt{commit\_material\_state()}
through virtual dispatch, enabling heterogeneous element containers.

\subsection{Layer 6: Analysis (Solver Boundary)}
\label{sec:layer6-analysis}

Files:
\texttt{src/analysis/NLAnalysis.hh},
\texttt{src/analysis/DynamicAnalysis.hh}

\subsubsection{NonlinearAnalysis (SNES)}

The SNES-driven Newton--Raphson solver provides two callbacks:
\begin{itemize}
  \item \texttt{FormResidual(u)} ---
        assembles $\mathbf{R} = \mathbf{f}_{\mathrm{int}}(\mathbf{u}) -
        \mathbf{f}_{\mathrm{ext}}$
        via element-level \texttt{compute\_internal\_force\_vector()}.
  \item \texttt{FormJacobian(u)} ---
        assembles $\mathbf{K}_t(\mathbf{u})$
        via element-level \texttt{compute\_tangent\_stiffness\_matrix()}.
\end{itemize}

Both callbacks are \textbf{side-effect free}: they evaluate stress
and tangent at the current displacement without modifying internal
variables.  The commit path is:

\begin{verbatim}
void commit_state() {
    // Scatter global -> local (+ imposed BCs)
    DMGlobalToLocal(dm, U, INSERT_VALUES, u_local);
    VecAXPY(u_local, 1.0, model_->imposed_solution());

    for (auto& element : model_->elements()) {
        element.commit_material_state(u_local);  // -> Layer 5 -> 4 -> 3 -> 2 -> 1
    }

    // Update model state for post-processing
    DMGlobalToLocal(dm, U, INSERT_VALUES, model_->state_vector());
}
\end{verbatim}

This private method is called \emph{only} when the SNES converged
reason is positive ($> 0$):

\begin{verbatim}
bool solve() {
    SNESSolve(snes_, nullptr, U);
    if (reason > 0) {
        commit_state();    // safe to commit
        return true;
    }
    // Diverged — do NOT commit corrupted state
    return false;
}
\end{verbatim}

\subsubsection{Incremental Loading with Bisection}

The \texttt{solve\_incremental(N, max\_bisections)} method applies
$N$ load steps with automatic bisection on divergence:

\begin{enumerate}
  \item Snapshot displacement $\mathbf{u}_{\text{snap}} \leftarrow \mathbf{u}$.
  \item Apply load level $\lambda \cdot \mathbf{f}_{\mathrm{ext}}$,
        call \texttt{SNESSolve}.
  \item If converged: call \texttt{commit\_state()}, proceed to
        next step.
  \item If diverged: revert $\mathbf{u} \leftarrow \mathbf{u}_{\text{snap}}$
        (material state unchanged because commit was never called),
        then bisect $[\lambda_{\text{prev}}, \lambda_{\text{target}}]$
        into two half-steps and retry recursively.
\end{enumerate}

This implements the key invariant:
\begin{quote}
  \emph{Material internal variables are never corrupted by a diverged
  step, because \texttt{commit()} is only called on convergence.}
\end{quote}

\subsubsection{DynamicAnalysis (PETSc TS)}

The \texttt{DynamicAnalysis} class uses PETSc TS for time integration
(Generalized-$\alpha$ or Newmark).  Material commit occurs in the TS
monitor callback after each converged time step, using the identical
element-level \texttt{commit\_material\_state()} path.


%% ------------------------------------------------------------------
\section{Audit Summary: Existing Constitutive Models}
\label{sec:phase1-audit-models}

\subsection{Elastic Path}

\begin{itemize}
  \item \texttt{ElasticRelation<Policy>} satisfies
        \texttt{ConstitutiveRelation}: $\boldsymbol{\sigma} =
        \mathbb{C} : \boldsymbol{\varepsilon}$, constant tangent.
  \item Injected via \texttt{PassThroughIntegrator} (elastic update).
  \item \texttt{commit()} is a no-op (path-independent).
  \item Verified by SNES tests with convergence in 1 iteration.
\end{itemize}

\subsection{$J_2$ Plasticity}

\begin{itemize}
  \item \texttt{PlasticityRelation<Policy, YieldF, Hardening, Flow,
        YieldFn, ReturnAlg>} implements composable classical plasticity.
  \item Components:
    \begin{itemize}
      \item \texttt{VonMises} yield function,
      \item \texttt{LinearIsotropicHardening} / mixed hardening,
      \item \texttt{AssociatedFlow} rule,
      \item \texttt{ScalarConsistencyReturn} algorithm (radial return).
    \end{itemize}
  \item Full return mapping:
        \begin{align}
          \boldsymbol{\varepsilon}^{\text{tr}}
            &= \boldsymbol{\varepsilon}_{n+1}
              - \boldsymbol{\varepsilon}^p_n, \\
          f_{\text{trial}}
            &= \|\text{dev}(\mathbb{C}:\boldsymbol{\varepsilon}^{\text{tr}})\|
              - \sigma_y(\alpha_n), \\
          \text{if } f_{\text{trial}} &\le 0:
            \quad \text{elastic step}, \\
          \text{else}: &\quad \text{solve consistency equation for}
            \;\Delta\gamma, \\
          \boldsymbol{\varepsilon}^p_{n+1}
            &= \boldsymbol{\varepsilon}^p_n
              + \Delta\gamma \, \mathbf{n},
            \qquad \alpha_{n+1} = \alpha_n + \Delta\gamma.
        \end{align}
  \item Internal variables: plastic strain
        $\boldsymbol{\varepsilon}^p$, equivalent plastic strain
        $\bar{\varepsilon}^p$, backstress $\boldsymbol{\beta}$
        (kinematic hardening).
  \item Exposes both external-state API
        (\texttt{compute\_response($\varepsilon$, $\alpha$)}) and
        internal-state API (\texttt{compute\_response($\varepsilon$)}).
  \item Injected via \texttt{EmbeddedRelationIntegrator} (inelastic
        update).
  \item Consistent tangent $\mathbb{C}^{ep}$ computed as:
        \[
          \mathbb{C}^{ep} = \mathbb{C}^e
            - \frac{2\mu\,\Delta\gamma}{\|\mathbf{s}^{\mathrm{tr}}\|}
              \left( \mathbb{C}^e - \tfrac{1}{3}\text{tr}(\mathbb{C}^e)
                \otimes \mathbf{I} \right)
            - 2\mu \left(\frac{1}{1+H/(3\mu)}
                         - \frac{\Delta\gamma}{\|\mathbf{s}^{\mathrm{tr}}\|/(2\mu)}
                   \right)
              \mathbf{n} \otimes \mathbf{n}.
        \]
  \item Tested in
        \texttt{test\_plasticity\_composition.cpp} (10 unit tests)
        and through SNES-level tests.
\end{itemize}

\subsection{Hyperelastic Models}

\begin{itemize}
  \item \texttt{SVKRelation}: Saint-Venant--Kirchhoff,
        $\Psi = \tfrac{\lambda}{2}(\text{tr}\,\mathbf{E})^2
                + \mu\,\mathbf{E}:\mathbf{E}$
  \item \texttt{NeoHookeanRelation}: compressible Neo-Hookean,
        $\Psi = \tfrac{\mu}{2}(I_1 - 3) - \mu\ln J
                + \tfrac{\lambda}{2}(\ln J)^2$
  \item Both support Total Lagrangian and Updated Lagrangian kinematic
        policies.  Tested incrementally under large deformation through
        SNES.
\end{itemize}

\subsection{Fiber Section Models}

\begin{itemize}
  \item \texttt{FiberSection} aggregates uniaxial fibers
        (steel + concrete).
  \item Uniaxial materials: \texttt{MenegottoPintoSteel},
        \texttt{KentParkConcrete} --- both have internal
        \texttt{commit()} that persists state.
  \item Commit chain: \texttt{BeamElement} $\to$ \texttt{MaterialSection}
        $\to$ \texttt{Material<BeamPolicy>} $\to$ \texttt{FiberSection}
        $\to$ per-fiber \texttt{commit()}.
  \item Tested in \texttt{test\_rc\_fiber\_frame\_nonlinear.cpp}
        (SNES-level).
\end{itemize}

\subsection{Finite Strain Damage}

\begin{itemize}
  \item \texttt{FiniteStrainDamageRelation} ---
        isotropic damage with energy-based criterion.
  \item Uses \texttt{ContinuumLocalProblemIntegrator} with
        Newton local solver for the internal-variable update.
  \item Tested in \texttt{test\_finite\_strain\_damage.cpp}.
\end{itemize}


%% ------------------------------------------------------------------
\section{Verification: End-to-End Commit Chain Test}
\label{sec:phase1-verification}

The existing test suite already exercises the commit chain
through several SNES integration tests:

\begin{itemize}
  \item \texttt{test\_snes\_hyperelastic.cpp} ---
        SmallStrain + TotalLagrangian + NeoHookean/SVK,
        incremental loading with convergence verification.
  \item \texttt{test\_plasticity\_composition.cpp} ---
        Unit-level tests for the composable plasticity architecture,
        including external-algorithmic-state path verification.
  \item \texttt{test\_rc\_fiber\_frame\_nonlinear.cpp} ---
        Full structural beam with fiber sections under SNES.
  \item \texttt{test\_finite\_strain\_damage.cpp} ---
        Continuum local-problem integrator with damage.
\end{itemize}

These tests verify that:
\begin{enumerate}
  \item SNES converges for elastic (1 iteration) and inelastic
        (multiple iterations) problems.
  \item Incremental loading with bisection correctly commits state
        only after convergence.
  \item The full chain from solver through type-erasure to
        constitutive relation produces correct force-displacement
        responses.
\end{enumerate}


%% ------------------------------------------------------------------
\section{Remaining Gaps for Future Phases}
\label{sec:phase1-gaps}

While the constitutive state protocol is functionally complete for the
SNES and PETSc~TS solvers, three gaps remain for later phases:

\begin{enumerate}
  \item \textbf{No \texttt{revert()} in type-erased interface.}
        \texttt{Material<Policy>} exposes \texttt{commit()} but not
        \texttt{revert()}.  The \texttt{revert\_trial()} method exists
        at the storage and state layers but is not wired through
        \texttt{MaterialConcept}'s virtual table.
        This is safe for SNES (implicit revert by not committing),
        but arc-length or substepping algorithms may need explicit
        revert at the element level.
        \emph{Planned for Phase~3.}

  \item \textbf{No \texttt{revert\_material\_state()} in
        \texttt{FiniteElement} concept.}
        The concept requires \texttt{commit\_material\_state(u)} but
        not a corresponding revert operation.  If the solver needs
        to undo a partially committed sub-step (e.g.\ in a load-arc
        method), this would need to be added.
        \emph{Planned for Phase~3.}

  \item \textbf{Commit logic duplicated in Analysis classes.}
        Both \texttt{NonlinearAnalysis} and \texttt{DynamicAnalysis}
        contain identical commit-state loops (scatter global$\to$local,
        loop elements, commit).  Extracting this to a Model-level
        \texttt{commit\_state(Vec u\_global)} method would reduce
        duplication and provide a single commit entry point.
        \emph{Planned for Phase~3.}
\end{enumerate}


%% ------------------------------------------------------------------
\section{Architectural Diagram: Commit Chain}
\label{sec:phase1-architecture-diagram}

The following stack trace illustrates the commit path from solver
to storage for a continuum element with inelastic material:

\begin{verbatim}
NonlinearAnalysis::commit_state()
 │  DMGlobalToLocal(dm, U, INSERT_VALUES, u_local)
 │
 └─► for each element:
      ContinuumElement::commit_material_state(u_local)
       │  u_e = extract_element_dofs(u_local)
       │  kin = KinematicPolicy::compute_kinematics(B, u_e, ...)
       │
       └─► MaterialPoint::commit(kin)
            │
            └─► Material<Policy>::commit(kin)    [virtual dispatch]
                 │
                 └─► OwningMaterialModel::commit(kin)
                      │
                      └─► EmbeddedRelationIntegrator::commit(site, k)
                           │  relation.update(k)
                           │    └─► PlasticityRelation: return mapping,
                           │        update ε_p, α, backstress
                           │
                           └─► site.constitutive_state().commit_trial()
                                │
                                └─► TrialCommittedState::commit_trial()
                                     committed_ ← trial_
                                     has_trial_ ← false
\end{verbatim}

For elastic materials, the same path reaches
\texttt{PassThroughIntegrator::commit()}, which is a no-op.

For structural elements (beams/shells), the path goes through
\texttt{MaterialSection::commit(strain)} instead of
\texttt{MaterialPoint::commit(kin)}, but reaches the same
underlying \texttt{Material<>::commit()} via the section's material.

%% ------------------------------------------------------------------
\section{Conclusions}
\label{sec:phase1-conclusions}

Phase~1 of the master plan is \textbf{complete}.  The constitutive
state protocol---trial evaluation, commit on convergence, implicit
revert on divergence---is implemented across all six architectural
layers and verified by the existing test suite.  The key findings are:

\begin{enumerate}
  \item The design documented in Chapters~28--31 was already implemented
        in code, resolving the design-implementation gap identified in
        Chapter~\ref{ch:convergence-master-plan} for this domain.
  \item The SNES solver correctly exploits the side-effect-free
        evaluation property, committing state only after verified
        convergence and automatically bisecting on divergence.
  \item The three remaining gaps (explicit revert at type-erased
        boundary, revert in \texttt{FiniteElement} concept, and
        commit duplication in Analysis classes) are non-blocking
        and are scheduled for Phase~3.
\end{enumerate}

With Phase~0 (RAII) and Phase~1 (constitutive protocol) both
verified, the next step is Phase~2: heterogeneous element containers
to support mixed continuum + structural assemblages.
